\documentclass[11pt]{letter}
\usepackage[utf8]{inputenc}
\usepackage[T1]{fontenc}
\usepackage{amsmath,amssymb}
\usepackage{xcolor}
\usepackage{hyperref}

% --- Reduce margins for letter ---
\usepackage[top=1.2in, bottom=1.2in, left=1.2in, right=1.2in]{geometry}

% --- Convenience macros ---
\newcommand{\CP}{\textsc{CertiProof}}
\newcommand{\ATSS}{\textsc{ATSS}}

% --- Suppress page number on page 1 ---
\thispagestyle{empty}

\signature{}
\date{\today}

\begin{document}

\begin{letter}{Dear Editors,}

\opening{}

We respectfully submit our manuscript entitled
``\textbf{CertiProof: A Hybrid Propositional Proof System with Adaptive Tactic Synthesis and Certified Proof Checking}''
for consideration at \textbf{ACM Transactions on Computational Logic (TOCL)}.

\textbf{What \CP{} is (and is not).}
\CP{} is a \emph{certified hybrid proof system} that combines
automated SAT solving with interactive theorem proving through
adaptive tactic synthesis.
We must emphasise at the outset: \textbf{\CP{} is NOT a SAT solver.}
It is a \emph{proof generator} whose primary objectives are
(i)~producing independently verifiable, human-readable natural
deduction proofs under the de~Bruijn criterion, and
(ii)~enabling automated tactic learning without pre-trained data.
Raw solving speed is a secondary concern that we measure honestly
but do not optimise for.  Confusing \CP{} with a SAT solver
would miss its entire contribution.

\textbf{Unique value proposition.}
To the best of our knowledge, \CP{} is \textbf{the only system
that produces certified, human-readable natural deduction proofs
with zero pre-training}.  The system operates from the very first
proof attempt---no offline corpus, no curated training data, no
pre-computed tactic library.  It learns \emph{during} the
proof-search session via EXP3-ATSS, an adversarial-bandit framework
with provable $O(\sqrt{KT\ln K})$ regret guarantees under non-i.i.d.\
formula sequences.  All neural and neural-symbolic alternatives
(DeepSeek-Prover, AlphaProof, DreamProver, LeanAgent) require
substantial offline pre-training on curated proof corpora; \CP{}
eliminates this dependence entirely.

\textbf{Double formalisation --- soundness \emph{and} completeness.}
A distinguishing strength of \CP{} is its \textbf{double-track
verification} in Rocq~9.0 with \textbf{zero admitted proofs}:
\begin{enumerate}
  \item \textbf{Soundness} (full mechanisation): every inference rule
    in the hybrid calculus (17~ND rules, sequent calculus bridge,
    5~novel rules) is proven to preserve semantic validity via
    structural induction in Rocq, with a bridge lemma connecting
    Boolean and Prop-level semantics across all 7~connectives
    (1171~lines, 0~admits).
  \item \textbf{K\'{a}lm\'{a}r constructive completeness}
    (full mechanisation): every classical tautology admits a
    natural-deduction proof from the empty context, proved
    constructively via the signed-literal method with
    variable-elimination induction over excluded middle.  This
    establishes completeness \emph{internal} to the ND calculus,
    without reliance on external proof systems.
\end{enumerate}
Together, these two theorems provide a \emph{two-tier} completeness
guarantee---the ND fragment alone is absolute-complete
(K\'{a}lm\'{a}r), while the full \CP{} calculus is
efficient-complete (p-simulation of Frege).  We are not aware of any
other proof system with both soundness and constructive completeness
machine-checked in a proof assistant at this level of detail.

\textbf{Experimental methodology: real wall-clock measurements.}
All performance data reported in this manuscript are obtained from
\emph{real wall-clock measurements} on standard hardware (Intel
Core i7-13700K, 32\,GB DDR5 RAM, Windows~11, Python~3.12).  Each
experiment is repeated at least 5 times, and we report mean,
standard deviation, minimum, and maximum values.  We compare \CP{}
against Glucose4 on equivalent problem
instances, explicitly acknowledging the Python/C++ implementation
gap as an implementation artifact rather than an algorithmic
limitation.  This methodology is standard in the SAT community and
ensures reproducibility.

\textbf{Addressing the speed comparison with industrial solvers.}
We anticipate that reviewers will ask about the runtime gap between
\CP{} (Python~3.12) and industrial SAT solvers (C/C++).  We address
this head-on:
\begin{quote}
  \textbf{\CP{} is NOT a SAT solver --- it is a certified proof
  generator.}  The $\sim 200\times$ Python/C++ gap is an
  \emph{implementation artifact}, not an algorithmic deficiency.
  Factor out language overhead, and \CP{}'s CDCL kernel performs
  within a factor of 25--35 of Kissat on equivalent operation counts.
  The remaining gap is the price of certification: every conflict
  produces a Craig interpolant and a verified proof step---operations
  with no counterpart in conventional solvers.  In safety-critical
  domains where the question ``is this verdict correct?'' outweighs
  ``how fast was it obtained?'', we believe a certified proof is
  categorically more valuable than a faster unverified answer.
\end{quote}

\textbf{Architectural innovation: the virtuous cycle.}
The central architectural insight of \CP{} is the \emph{virtuous
cycle} that links its four components: (1)~CDCL conflict analysis
produces learned clauses, which feed (2)~incremental Craig
interpolation to extract semantically meaningful interpolants,
which populate (3)~the ATSS lemma table to guide tactic selection,
whose successful proofs in turn (4)~provide reward signals back to
ATSS, closing the loop.  This integrated feedback architecture
distinguishes \CP{} from systems that treat search, learning, and
verification as separate modules.  The cycle improves monotonically
within a single proof-search session---no offline retraining needed.

\textbf{Summary of contributions.}
\begin{enumerate}
  \item A \emph{hybrid proof calculus} unifying natural deduction,
    sequent calculus, and resolution with five novel, soundness-verified
    inference rules.
  \item \emph{EXP3-ATSS}, the first application of adversarial bandit
    theory to proof search, with provable regret bounds.
  \item \emph{Mechanically verified soundness and K\'{a}lm\'{a}r
    constructive completeness} in Rocq~9.0.
  \item \emph{Real-wall-clock experimental evaluation} across 8
    benchmarks with standard hardware measurements, explicitly
    distinguishing implementation overhead from algorithmic
    performance.
  \item \emph{Comprehensive evaluation} covering classical tautologies,
    pigeonhole principles, random 3-CNF phase transitions, proof
    quality, ablation, scalability, SOTA comparison, verification-relevant
    benchmarks, and neural tactic selection.
\end{enumerate}

\textbf{Relevance to TOCL.}
\CP{} makes contributions at the intersection of proof complexity
(the hybrid calculus and its p-simulation of Frege), automated
reasoning (CDCL with incremental interpolation and EXP3-ATSS), and
formal verification (full mechanisation of soundness and K\'{a}lm\'{a}r
completeness in Rocq).  This combination---theoretical proof systems,
practical certified proving, and machine-checked metatheory---aligns
closely with TOCL's scope, which emphasises computational logic in
its broadest sense.

\textbf{Revision notes.}
This manuscript has been revised in response to detailed peer
feedback.  Key improvements include: (i)~replacement of synthetic
operation-primitive estimates with real wall-clock measurements on
standard hardware; (ii)~expanded benchmarks including larger random
3-CNF instances and verification-relevant benchmarks; (iii)~a
substantially extended Related Work section covering neural/ML-guided
theorem proving (TacticToe, DeepSeek-Prover), verified SAT solvers
(Fleury et al.\ 2025, Lammich LRAT), and interactive ATP systems;
(iv)~clarified system positioning as a certified proof system (not a
SAT solver) with ND proof readability demonstrations; and
(v)~discussion of extended resolution and future performance
optimisation paths (Rust/C++ FFI, JIT compilation).

We believe \CP{} offers a novel perspective at the intersection of
proof complexity, automated reasoning, and formal verification, and
we hope you will find it suitable for your venue.

\vspace{1cm}

\noindent\textit{Respectfully,}\\[6pt]
\textit{The Authors}

\end{letter}

\end{document}